% =====================================================================
%  ANEXO D. Registro de incidentes
% =====================================================================
\section{Registro de incidentes}
\label{anx:incidentes}

\begin{notanormativa}[Registro nuevo respecto de la plantilla original]
Corresponde a la actividad \textbf{IR2 «Crear o actualizar el reporte de incidente»} de
\normap{2}. Las reglas de análisis previo, las escalas de severidad y prioridad y el flujo de
estados se definen en el apartado~\ref{sec:gestion-incidentes}; aquí se registran los
incidentes concretos.
\end{notanormativa}

\begin{instruccion}
Reproduzca la ficha por cada incidente reportado. Los pasos de reproducción son el campo que
determina si el incidente podrá ser atendido: escríbalos de forma que una persona ajena al
equipo pueda reproducir el fallo sin preguntar nada.

Recuerde que \textbf{severidad y prioridad son campos distintos} y que un incidente puede tener
como origen el entorno, los datos o el propio caso de prueba, y no el objeto probado.
\end{instruccion}

\subsection*{Ficha de incidente}

\begin{xltabular}{\textwidth}{|E{4.4cm}|Y|}
\caption{Formato de registro de un incidente de prueba.}\label{tab:ficha-incidente}\\ \hline
\endfirsthead
\multicolumn{2}{@{}l@{}}{\footnotesize\itshape\tablename~\thetable{} (continuación)}\\ \hline
\endhead
\multicolumn{2}{@{}r@{}}{\footnotesize\itshape continúa en la página siguiente}\\
\endfoot
\endlastfoot
  Identificador                & \campo{INC-05} \\ \hline
  Resumen                      & \campo{Una línea que identifique el problema} \\ \hline
  Descripción                  & \campo{Explicación detallada de lo observado} \\ \hline
  Origen determinado (IR1)     & \campo{Objeto de prueba / entorno / datos / caso de prueba / duplicado} \\ \hline
  Ejecución que lo reveló      & \campo{EJ-001 (\ref{anx:ejecuciones})} \\ \hline
  Objeto de prueba afectado    & \campo{OP-01} \\ \hline
  Resultado esperado           & \campo{Qué debía ocurrir} \\ \hline
  Resultado obtenido           & \campo{Qué ocurrió} \\ \hline
  Pasos de reproducción        & \campo{1. …\quad 2. …\quad 3. …} \\ \hline
  ¿Reproducible?               & \campo{Siempre / intermitente / no reproducido} \\ \hline
  Versión del software         & \campo{v1.2.0} \\ \hline
  Entorno y configuración      & \campo{ENT-xx y parámetros} \\ \hline
  Datos utilizados             & \campo{DAT-xx} \\ \hline
  Evidencia                    & \campo{Capturas, registros, archivos} \\ \hline
  Severidad                    & \campo{Cuadro~\ref{tab:severidad}} \\ \hline
  Prioridad                    & \campo{Cuadro~\ref{tab:prioridad}} \\ \hline
  Riesgo de producto asociado  & \campo{RP-01} \\ \hline
  Reportado por / fecha        & \campo{Nombre — dd/mm/aaaa} \\ \hline
  Responsable asignado         & \campo{Nombre o equipo} \\ \hline
  Estado                       & \campo{Cuadro~\ref{tab:estados-incidente}} \\ \hline
  Fecha de última actualización& \campo{dd/mm/aaaa} \\ \hline
  Resolución o disposición     & \campo{Qué se hizo, o por qué se difiere o rechaza} \\ \hline
  Ejecución de confirmación    & \campo{EJ-014, si ya se confirmó la corrección} \\ \hline
\end{xltabular}

El Cuadro~\ref{tab:ficha-incidente} contiene los campos mínimos para que un incidente pueda
comunicarse, priorizarse y cerrarse de forma verificable. El campo de origen evita que un
problema del entorno se contabilice como defecto del producto.

\subsection*{Registro consolidado de incidentes}

\begin{longtable}{|C{1.4cm}|L{3.6cm}|C{1.5cm}|C{1.5cm}|C{1.6cm}|C{1.8cm}|C{1.6cm}|}
\caption{Registro consolidado de incidentes.}
\label{tab:registro-incidentes} \\
\hline
\filaenc \textbf{ID} & \textbf{Resumen} & \textbf{Sever.} & \textbf{Prior.} & \textbf{Estado} & \textbf{Responsable} & \textbf{Ejecución} \\ \hline
\endfirsthead
\hline
\filaenc \textbf{ID} & \textbf{Resumen} & \textbf{Sever.} & \textbf{Prior.} & \textbf{Estado} & \textbf{Responsable} & \textbf{Ejecución} \\ \hline
\endhead
\hline \multicolumn{7}{|r|}{\footnotesize\textit{continúa en la página siguiente}} \\ \hline
\endfoot
\hline
\endlastfoot
\campo{INC-05} & \campo{Resumen} & \campo{Alta} & \campo{Normal} & \campo{Cerrado} & \campo{Nombre} & \campo{EJ-001} \\ \hline
 & & & & & & \\ \hline
 & & & & & & \\ \hline
\end{longtable}

El Cuadro~\ref{tab:registro-incidentes} es la vista de conjunto que alimenta el seguimiento
(apartado~\ref{sec:seguimiento}) y, al cierre, el
Cuadro~\ref{tab:incidentes-abiertos}.
